\documentclass[12pt]{article}

\usepackage[utf8]{inputenc}
\usepackage[T1]{fontenc}
\usepackage[brazil]{babel}
\usepackage{geometry}
\usepackage{graphicx}
\usepackage{url}
\usepackage{xurl}
\usepackage[colorlinks=true, allcolors=black]{hyperref}
\usepackage{amsmath}
\usepackage{booktabs}
\usepackage{times}
\usepackage{microtype}
\usepackage{siunitx}

\geometry{
  a4paper,
  top=3.5cm,
  bottom=2.5cm,
  left=3.0cm,
  right=3.0cm
}

\linespread{1.0}
\setlength{\parindent}{1.27cm}
\setlength{\parskip}{6pt}

\sisetup{
  output-decimal-marker = {,},
  group-separator       = {.},
  group-minimum-digits  = 4
}

\renewenvironment{abstract}{
  \small
  \begin{list}{}{
    \setlength{\leftmargin}{0.8cm}
    \setlength{\rightmargin}{0.8cm}
  }
  \item[\hskip \labelsep {\bfseries Abstract.}]
}{
  \end{list}
}

\newenvironment{resumo}{
  \small
  \begin{list}{}{
    \setlength{\leftmargin}{0.8cm}
    \setlength{\rightmargin}{0.8cm}
  }
  \item[\hskip \labelsep {\bfseries Resumo.}]
}{
  \end{list}
}

\title{\bfseries Análise Comparativa de Desempenho de Algoritmos de Ordenação
       em Diferentes Cenários de Entrada}

\author{Felipe Rochoel Costa}

\date{}

\begin{document}

\maketitle

\begin{center}
  \small
  Graduação em Engenharia de Computação \\
  Universidade Tecnológica Federal do Paraná (UTFPR) -- Campus Apucarana\\
  Apucarana -- PR -- Brasil \\
  \texttt{feliperochoelcosta@alunos.utfpr.edu.br}
\end{center}

\vspace{0.5cm}

\begin{abstract}
  This paper presents an empirical and comparative performance analysis of
  three widely studied sorting algorithms: Bubble Sort, Merge Sort, and
  Quick Sort. The evaluation focuses on execution time and the number of key
  comparisons and element swaps across three initial data ordering scenarios
  (random, inversely sorted, and nearly sorted lists) with dataset sizes
  ranging from 10{,}000 to 1{,}000{,}000 elements. The experimental results
  validate the theoretical asymptotic complexities, highlight the overhead of
  quadratic algorithms in large volumes of data, and explore the practical
  efficiency trade-offs between divide-and-conquer approaches. Source code
  and experimental datasets are publicly available at
  \url{https://github.com/rochoel-felipe/analise-algoritmos-ordenacao}.
\end{abstract}

\begin{resumo}
  Este artigo apresenta uma análise empírica e comparativa de desempenho de
  três algoritmos de ordenação amplamente estudados: Bubble Sort, Merge Sort
  e Quick Sort. A avaliação centra-se no tempo de execução e no número de
  comparações de chaves e trocas de elementos sob três cenários de ordenação
  inicial dos dados (listas aleatórias, inversamente ordenadas e quase
  ordenadas), com tamanhos de conjuntos de dados variando de 10.000 a
  1.000.000 de elementos. Os resultados experimentais validam as
  complexidades assintóticas teóricas, evidenciam o gargalo de algoritmos
  quadráticos em grandes volumes de dados e exploram as compensações de
  eficiência prática entre as abordagens de divisão e conquista. O
  código-fonte e os dados experimentais estão disponíveis publicamente em
  \url{https://github.com/rochoel-felipe/analise-algoritmos-ordenacao}.
\end{resumo}

% ===================================================================
\section{Introdução}
% ===================================================================

A ordenação de dados constitui um dos problemas mais fundamentais e
exaustivamente explorados no âmbito da Ciência da Computação e da Engenharia
de Software. Definida como o processo de arranjar uma coleção de registros em
uma sequência linear determinada por uma relação de ordem total, a ordenação
eficiente é um requisito crítico para a otimização de operações subsequentes,
tais como buscas binárias, eliminação de duplicatas e processamento de bancos
de dados.

Com o advento do \textit{Big Data} e o crescimento exponencial do volume de
informações manipuladas por sistemas computacionais modernos, a escolha do
algoritmo de ordenação ideal deixou de ser uma mera decisão conceitual para
tornar-se um fator determinante no consumo de recursos, na latência e na
escalabilidade de softwares. Diferentes algoritmos apresentam comportamentos
distintos que dependem diretamente não apenas do número de elementos a
ordenar, mas também do estado de ordenação prévio desse conjunto de dados.

Este trabalho tem como objetivo central realizar um estudo comparativo e
experimental de três algoritmos de ordenação proeminentes: Bubble Sort, Merge
Sort e Quick Sort. A análise avalia o impacto do tamanho da entrada (vetores
de $n \in \{10.000;\, 100.000;\, 1.000.000\}$ elementos) e o comportamento
sob diferentes distribuições iniciais dos dados (listas aleatórias,
inversamente ordenadas e quase ordenadas). Por meio da mensuração do tempo
computacional de execução e da contagem de operações primitivas (comparações
e trocas), busca-se confrontar os achados práticos com as predições
matemáticas da análise assintótica clássica.

% ===================================================================
\section{Fundamentação Teórica}
% ===================================================================

A eficiência de um algoritmo é formalmente avaliada por meio da notação
assintótica grande-$O$, que descreve o limite superior do crescimento do
tempo de execução em função do tamanho da entrada~$n$. Os algoritmos
selecionados pertencem a duas classes principais de complexidade e estratégias
de projeto distintas~\cite{cormen,knuth}.

\subsection{O Algoritmo Bubble Sort e sua Complexidade}

O Bubble Sort é um algoritmo elementar iterativo e baseado em comparações. O
princípio operacional consiste em percorrer o vetor repetidamente a partir do
início, comparando pares de elementos adjacentes e realizando uma operação de
troca (\textit{swap}) sempre que dois elementos consecutivos estiverem fora da
ordem desejada. Esse procedimento faz com que os maiores elementos
``flutuem'' gradualmente para as posições finais.

Do ponto de vista analítico, o Bubble Sort possui complexidade de tempo de
pior caso e caso médio em $O(n^2)$. O pior caso ocorre quando o vetor está
em ordem estritamente inversa, exigindo $n(n-1)/2$ comparações e o mesmo
número de trocas. Com a otimização de uma variável sinalizadora
(\textit{flag}) que interrompe a execução caso nenhuma troca ocorra em uma
passagem completa, o algoritmo atinge $O(n)$ no melhor caso — quando o vetor
já está ordenado. Sua complexidade de espaço é $O(1)$, classificando-o como
algoritmo \textit{in-place}. Apesar de didaticamente difundido e simples de
implementar, sua eficiência degrada de maneira proibitiva em conjuntos de
dados massivos.

\subsection{O Algoritmo Merge Sort e sua Complexidade}

O Merge Sort é um algoritmo de ordenação estável fundamentado no paradigma de
``dividir para conquistar'', proposto por John von Neumann em
1945~\cite{knuth}. O funcionamento divide-se em três etapas recursivas:

\begin{enumerate}
    \item \textbf{Divisão:} O vetor original é segmentado em dois subvetores
          de tamanho aproximadamente igual, $\lfloor n/2 \rfloor$ e
          $\lceil n/2 \rceil$.
    \item \textbf{Conquista:} Cada subvetor é ordenado recursivamente por
          novas chamadas ao Merge Sort. Quando um subvetor atinge tamanho
          unitário, está trivialmente ordenado.
    \item \textbf{Combinação (\textit{merge}):} Os subvetores ordenados são
          intercalados de forma linear para recompor um único vetor totalmente
          ordenado.
\end{enumerate}

A propriedade mais notável do Merge Sort é sua consistência temporal: garante
$O(n \log n)$ para melhor, pior e caso médio, pois a árvore de recorrência
mantém profundidade $\log n$ e o custo de intercalação em cada nível é $O(n)$.
No entanto, a operação de intercalação convencional não é \textit{in-place},
requerendo um vetor auxiliar de tamanho $O(n)$, o que pode limitar sua
viabilidade em ambientes com restrição severa de memória.

\subsection{O Algoritmo Quick Sort e sua Complexidade}

O Quick Sort, idealizado por C.~A.~R.~Hoare em 1959~\cite{knuth}, emprega o
paradigma de divisão e conquista em torno de um elemento de referência
denominado \textit{pivô}. A operação crítica é o particionamento, que
rearranja os elementos de modo que todos os itens menores que o pivô fiquem à
sua esquerda e os maiores ou iguais, à sua direita. Após o particionamento, o
pivô ocupa sua posição definitiva e o processo é aplicado recursivamente nos
dois subvetores resultantes.

Em cenários médios e favoráveis, o Quick Sort exibe complexidade $O(n \log n)$
e, devido à constante multiplicativa baixa e à excelente localidade de
referência para caches de CPU, costuma superar o Merge Sort na prática.
Contudo, sua eficiência é altamente sensível à escolha do pivô: quando o
pivô é sistematicamente o menor ou o maior elemento do conjunto — situação
comum em entradas já ordenadas ou inversamente ordenadas com pivô fixo no
primeiro elemento —, a árvore de chamadas recursivas atinge profundidade~$n$,
degradando a complexidade para $O(n^2)$. O consumo de espaço na pilha de
execução varia de $O(\log n)$ no caso médio a $O(n)$ no pior cenário.

% ===================================================================
\section{Método e Procedimentos}
% ===================================================================

\subsection{Ambiente de Testes}

Os algoritmos foram implementados em linguagem C e compilados com o GCC
(MinGW). Os experimentos foram executados de forma isolada na seguinte
configuração de hardware e software:

\begin{itemize}
    \item \textbf{Processador:} AMD Ryzen 5 5600G (3,90\,GHz)
    \item \textbf{Memória RAM:} 16\,GB DDR4 (com pilha de 64\,MB alocada
          dinamicamente para o Quick Sort)
    \item \textbf{Sistema Operacional:} Windows 11 Pro (64\,bits, versão 24H2)
    \item \textbf{Compilador:} GCC (MinGW) via Visual Studio Code
\end{itemize}

A medição do tempo de execução foi realizada com temporizadores de alta
precisão da biblioteca \texttt{<time.h>}, capturando o intervalo exato entre
o início e o término da rotina de ordenação. Para quantificar o esforço
computacional de forma independente do hardware, contadores internos foram
inseridos no código-fonte para registrar o número exato de comparações de
chaves realizadas e de trocas de posição executadas.

O código-fonte, os roteiros de testes e os conjuntos de dados brutos estão
disponibilizados publicamente no repositório:
{\small\url{https://github.com/rochoel-felipe/analise-algoritmos-ordenacao}}.

\subsection{Cenários Avaliados}

Cada algoritmo foi submetido a testes combinando três tamanhos de vetores
($n = 10.000$, $n = 100.000$ e $n = 1.000.000$) em três perfis de ordenação
inicial. Uma restrição laboratorial foi aplicada ao Bubble Sort, limitando
seu teste máximo a $n = 500.000$ para evitar sobrecarga temporal inviável;
para $n = 1.000.000$, o tempo esperado foi extrapolado a partir dos dados
coletados. Os perfis definidos foram:

\begin{itemize}
    \item \textbf{Listas Aleatórias:} Elementos gerados pseudoaleatoriamente
          com semente fixa (\texttt{srand(42)}), garantindo reprodutibilidade
          entre as rodadas.
    \item \textbf{Listas Inversamente Ordenadas:} Vetores preenchidos em
          ordem estritamente decrescente — pior caso teórico tanto para o
          Bubble Sort quanto para o Quick Sort com pivô fixo.
    \item \textbf{Listas Quase Ordenadas:} Vetores originalmente ordenados de
          forma crescente nos quais aproximadamente 2\% dos elementos tiveram
          suas posições deliberadamente permutadas.
\end{itemize}

Cada configuração foi executada em três rodadas independentes. As métricas de
comparações e trocas são determinísticas (idênticas entre rodadas); os tempos
de processamento reportados representam a média aritmética das três execuções.

% ===================================================================
\section{Resultados e Discussão}
% ===================================================================

\subsection{Desempenho em Listas Aleatórias}

A Tabela~\ref{tab:aleatorias} apresenta o comportamento dos algoritmos sobre
dados gerados arbitrariamente.

\begin{table}[ht]
\centering
\caption{Desempenho dos algoritmos de ordenação em listas aleatórias.}
\label{tab:aleatorias}
\begin{tabular}{l r r r r}
\toprule
\textbf{Algoritmo} & \textbf{Tamanho ($n$)} &
  \textbf{Tempo Médio (ms)} & \textbf{Comparações} & \textbf{Trocas} \\
\midrule
Bubble Sort & 10.000    & 126,67     & 49.994.535       & 24.845.189      \\
            & 100.000   & 21.494,33  & 499.844.889      & 250.489.970     \\
            & 500.000\textsuperscript{*}
                        & 544.215,00 & 124.999.422.355  & 62.425.097.729  \\
\midrule
Merge Sort  & 10.000    & 2,67       & 120.439          & 133.616         \\
            & 100.000   & 21,33      & 1.536.713        & 1.668.928       \\
            & 1.000.000 & 224,33     & 18.674.332       & 19.951.424      \\
\midrule
Quick Sort  & 10.000    & $<$1       & 166.685          & 78.147          \\
            & 100.000   & 8,33       & 2.090.695        & 1.001.984       \\
            & 1.000.000 & 100,00     & 35.550.416       & 10.268.671      \\
\bottomrule
\end{tabular}
\begin{flushleft}
\footnotesize \textsuperscript{*}Limitado a 500.000 elementos em razão da
complexidade quadrática do algoritmo.
\end{flushleft}
\end{table}

No cenário aleatório, os resultados práticos alinham-se com os limites
teóricos da análise assintótica. O crescimento do tempo e das operações do
Bubble Sort evidencia o custo quadrático $O(n^2)$: ao escalar de 10.000 para
100.000 elementos (fator 10), o tempo médio passou de 126,67\,ms para
21.494,33\,ms (aumento de aproximadamente 170 vezes, próximo da escala
$10^2$). Para $n = 500.000$, o tempo ultrapassou 9 minutos.

Os algoritmos baseados em divisão e conquista comprovaram sua eficiência
superior. O Merge Sort manteve-se linearitimicamente estável, ordenando
1.000.000 de registros em 224,33\,ms. O Quick Sort consolidou-se como o mais
veloz neste cenário: com tempo inferior a 1\,ms para $n = 10.000$ e apenas
100,00\,ms para o vetor máximo. Essa vantagem sobre o Merge Sort decorre da
menor constante multiplicativa e da ausência de alocação de memória auxiliar.

\subsection{Desempenho em Listas Inversamente Ordenadas}

Os dados para o cenário de ordenação inversa constam na
Tabela~\ref{tab:inversas}.

\begin{table}[ht]
\centering
\caption{Desempenho dos algoritmos em listas ordenadas de forma inversa.}
\label{tab:inversas}
\begin{tabular}{l r r r r}
\toprule
\textbf{Algoritmo} & \textbf{Tamanho ($n$)} &
  \textbf{Tempo Médio (ms)} & \textbf{Comparações} & \textbf{Trocas} \\
\midrule
Bubble Sort & 10.000    & 152,00      & 49.995.000        & 49.995.000        \\
            & 100.000   & 18.448,33   & 4.999.950.000     & 4.999.950.000     \\
            & 500.000\textsuperscript{*}
                        & 459.820,67  & 124.999.750.000   & 124.999.750.000   \\
\midrule
Merge Sort  & 10.000    & $<$1        & 64.608            & 133.616           \\
            & 100.000   & 19,67       & 815.024           & 1.668.928         \\
            & 1.000.000 & 158,33      & 9.884.992         & 19.951.424        \\
\midrule
Quick Sort  & 10.000    & 106,33      & 49.995.000        & 25.009.999        \\
            & 100.000   & 10.002,00   & 4.999.950.000     & 2.500.099.999     \\
            & 1.000.000 & 1.000.239,67& 499.999.500.000   & 250.000.999.999   \\
\bottomrule
\end{tabular}
\begin{flushleft}
\footnotesize \textsuperscript{*}Limitado a 500.000 elementos em razão da
complexidade quadrática do algoritmo.
\end{flushleft}
\end{table}

O cenário inversamente ordenado revelou comportamentos de extrema relevância
teórica. A escolha estática do primeiro elemento como pivô forçou o Quick
Sort a incorrer sistematicamente em particionamentos degenerados e
completamente desbalanceados (subvetores de tamanhos $0$ e $n-1$), degradando
sua complexidade para o pior caso $O(n^2)$. Para $n = 1.000.000$, o
algoritmo realizou quase 500 bilhões de comparações, demandando
1.000.239,67\,ms — aproximadamente 16,7 minutos — de processamento.

O Bubble Sort também atingiu seu ápice de ineficiência: cada comparação
resultou em uma troca, gerando a quantidade máxima de operações
$n(n-1)/2$. O Merge Sort, por outro lado, demonstrou resiliência assintótica
completa: indiferente à disposição desfavorável dos dados, manteve
158,33\,ms para o vetor máximo, consagrando-se como a abordagem mais segura
e robusta diante de distribuições inversas.

\subsection{Desempenho em Listas Quase Ordenadas}

Os dados para vetores quase totalmente ordenados constam na
Tabela~\ref{tab:quase}.

\begin{table}[ht]
\centering
\caption{Desempenho dos algoritmos em listas quase ordenadas.}
\label{tab:quase}
\begin{tabular}{l r r r r}
\toprule
\textbf{Algoritmo} & \textbf{Tamanho ($n$)} &
  \textbf{Tempo Médio (ms)} & \textbf{Comparações} & \textbf{Trocas} \\
\midrule
Bubble Sort & 10.000    & 76,33       & 49.806.809      & 1.286.410       \\
            & 100.000   & 4.019,33    & 2.723.276.019   & 39.529.560      \\
            & 500.000\textsuperscript{*}
                        & 22.727,67   & 15.617.581.497  & 144.816.562     \\
\midrule
Merge Sort  & 10.000    & 2,00        & 105.886         & 133.616         \\
            & 100.000   & 19,33       & 1.017.935       & 1.668.928       \\
            & 1.000.000 & 165,00      & 10.263.537      & 19.951.424      \\
\midrule
Quick Sort  & 10.000    & 1,67        & 2.021.659       & 149.163         \\
            & 100.000   & 2.305,00    & 226.830.634     & 537.932         \\
            & 1.000.000 & 477.187,33  & 467.783.498.400 & 1.280.759       \\
\bottomrule
\end{tabular}
\begin{flushleft}
\footnotesize \textsuperscript{*}Limitado a 500.000 elementos em razão da
complexidade quadrática do algoritmo.
\end{flushleft}
\end{table}

O cenário quase ordenado evidenciou dois comportamentos de grande interesse
prático. O Bubble Sort beneficiou-se expressivamente da pré-ordenação: com
apenas 144.816.562 trocas para $n = 500.000$ (contra 124.999.750.000 no
cenário inverso), e um tempo médio de 22.727,67\,ms, a redução é drástica em
relação ao pior caso. A otimização por \textit{flag} permite ao algoritmo
encerrar cada passagem mais cedo, diminuindo o número efetivo de iterações.
Ainda assim, o volume de comparações permaneceu elevado (cerca de 15 bilhões
para $n = 500.000$), pois o algoritmo precisa percorrer o vetor até confirmar
a ausência total de inversões.

O Quick Sort com pivô fixo no primeiro elemento também sofreu penalização
neste cenário, embora de forma menos severa que no caso inverso. Como os
primeiros elementos já se encontravam próximos de suas posições definitivas
na sequência crescente, as divisões geradas foram desbalanceadas, resultando
em 477.187,33\,ms para $n = 10^6$ e em curvas de crescimento próximas a
$O(n^2)$. Vale notar, contudo, que o número de trocas foi extremamente baixo
(apenas 1.280.759 para $n = 1.000.000$), confirmando que o custo recai quase
inteiramente sobre o particionamento, não sobre os rearranjos de elementos.
Tal comportamento reitera a necessidade prática de estratégias como pivô
aleatório ou mediana de três em implementações reais.

O Merge Sort reiterou sua estabilidade insensível ao arranjo inicial,
resolvendo o vetor máximo em 165,00\,ms — o menor tempo registrado entre
todos os algoritmos neste cenário.

% ===================================================================
\section{Conclusão}
% ===================================================================

Este estudo confrontou as formulações da teoria assintótica com o
comportamento empírico observado nos experimentos. O Bubble Sort provou-se
inadequado para conjuntos extensos — o ônus do crescimento quadrático
$O(n^2)$ é intransponível —, exibindo melhora expressiva apenas em
estruturas quase ordenadas, onde a otimização por \textit{flag} reduz
sensivelmente o número de trocas e, consequentemente, o tempo total.

As técnicas de divisão e conquista reafirmaram sua supremacia. O Merge Sort
destacou-se pela previsibilidade temporal absoluta, mantendo $O(n \log n)$
em todos os cenários ao custo de $O(n)$ de espaço auxiliar. O Quick Sort
consolidou-se como a ferramenta mais veloz no caso médio (dados aleatórios),
mas sua dependência crítica da estratégia de escolha do pivô foi empiricamente
exposta nos cenários inverso e quase ordenado, onde regrediu ao desempenho
quadrático. Notavelmente, neste último cenário, o Quick Sort realizou um
número mínimo de trocas, evidenciando que o gargalo reside exclusivamente na
fase de particionamento.

Conclui-se que a seleção do algoritmo ideal deve equilibrar não apenas o
volume de dados, mas o conhecimento prévio sobre a distribuição estrutural
da entrada. Para cenários de uso geral, o Merge Sort oferece maior segurança
e previsibilidade; para dados tipicamente aleatórios onde desempenho médio é
prioritário, o Quick Sort com escolha de pivô robusta (e.g., mediana de três
ou pivô aleatório) apresenta vantagens práticas relevantes.

% ===================================================================
\begin{thebibliography}{99}
% ===================================================================

\bibitem{cormen}
Cormen, T.~H., Leiserson, C.~E., Rivest, R.~L., and Stein, C. (2012).
\newblock {\em Algoritmos: Teoria e Prática}.
\newblock Elsevier, Rio de Janeiro, 3\textsuperscript{a}~edição.

\bibitem{slidesgoes}
Góes, W. (2026).
\newblock {\em Material Didático da Disciplina de Estruturas de Dados 2:
  Aulas de Bubble Sort, Merge Sort e Quick Sort}.
\newblock UTFPR, Campus Apucarana.

\bibitem{knuth}
Knuth, D.~E. (1998).
\newblock {\em The Art of Computer Programming, Volume~3: Sorting and
  Searching}.
\newblock Addison-Wesley, Reading, MA, 2nd~edition.

\end{thebibliography}

\end{document}
